\section{The five-row hand classification}\label{app:fivecases}

This appendix proves Proposition~\ref{prop:fivecases}. Throughout, $S^{\infty}$ is the
blackout of the five-row half-strip of all columns $x\le21$ defined in
Section~\ref{sec:five}. Row $1$ is black in every column $x\le21$ except $9$, $13$, $17$
and $21$. The black points in the other rows are
\[
A_2=\{1,9,17\},\qquad A_3=\{9,13\},\qquad A_4=\{5,13,21\},\qquad A_5=\{17,21\},
\]
where $A_j$ is the set of columns black in row $j$. Every column in $A_2\cup A_3\cup A_4\cup
A_5$ is congruent to $1$ modulo $4$, and every double column joins an odd row to an even row.
We call a column \emph{ordinary} if its only black point is in row $1$, and we write $B_x$ for
the set of black rows of column $x$ as in Section~\ref{sec:conv}. The set $F$ of black row
pairs is
\[
F=\bigl\{\{1,2\},\{1,4\},\{2,3\},\{3,4\},\{2,5\},\{4,5\}\bigr\},
\]
and in particular $\{1,3\}\notin F$ and $\{1,5\}\notin F$. A weak rectangle shows at most two
corners.

\begin{proposition}[the weak tilted rectangles of the half-strip]\label{prop:fivelist}
The tilted rectangles with corners in the half-strip fall into $18$ templates up to
horizontal translation. Only the $12$ templates that touch row $1$ have weak translates, and
these are the $59$ rectangles of Table~\ref{tab:fiveweak}, all with columns in $[-3,21]$,
$47$ showing two corners and $12$ showing one. Their presentations are pairwise distinct, and no
axis-aligned rectangle with corners in the half-strip has any of them as its presentation.
\end{proposition}

\begin{proof}
\emph{Step 1. The eighteen templates.} By Lemma~\ref{lem:leftmost}, a tilted rectangle has
a unique leftmost corner $(x,r)$, from which its two sides are $(a,-b)$ and $(c,d)$ with
positive integers satisfying $ac=bd$. Its corners are
$(x,r)$, $(x+a,r-b)$, $(x+c,r+d)$ and $(x+a+c,r+d-b)$, so it needs $r-b\ge1$ and $r+d\le5$,
that is, $b+d\le4$ and $b+1\le r\le5-d$. The pairs $(b,d)$ with $b+d\le4$ are $(1,1)$,
$(1,2)$, $(2,1)$, $(1,3)$, $(3,1)$ and $(2,2)$. Their products $bd$ have $1$, $2$, $2$, $2$,
$2$ and $3$ ordered factor pairs $(a,c)$, and the row $r$ has $3$, $2$, $2$, $1$, $1$ and $1$
choices. So there are $1\cdot3+2\cdot2+2\cdot2+2\cdot1+2\cdot1+3\cdot1=18$ templates, each
determined by $(a,b,c,d,r)$, and a tilted rectangle is a template together with its leftmost
column $x$. A template touches row $1$ exactly when $r=b+1$. The twelve templates with $r=b+1$
are listed in Table~\ref{tab:fivetemplates} with their corners. The six others are the unit
diamond with $r=3$ or $r=4$, the two templates $S_1$ and $S_2$ with $r=3$, and the two
templates $S_3$ and $S_4$ with $r=4$.

\begin{table}[t]
\centering\small
\begin{tabular}{@{}llll@{}}
\toprule
name & $(a,b,c,d)$ & corners & span \\
\midrule
$D$ & $(1,1,1,1)$ & $(x,2),(x+1,1),(x+1,3),(x+2,2)$ & $2$ \\
$S_1$ & $(1,1,2,2)$ & $(x,2),(x+1,1),(x+2,4),(x+3,3)$ & $3$ \\
$S_2$ & $(2,1,1,2)$ & $(x,2),(x+1,4),(x+2,1),(x+3,3)$ & $3$ \\
$S_3$ & $(1,2,2,1)$ & $(x,3),(x+1,1),(x+2,4),(x+3,2)$ & $3$ \\
$S_4$ & $(2,2,1,1)$ & $(x,3),(x+1,4),(x+2,1),(x+3,2)$ & $3$ \\
$Q_1$ & $(1,1,3,3)$ & $(x,2),(x+1,1),(x+3,5),(x+4,4)$ & $4$ \\
$Q_2$ & $(3,1,1,3)$ & $(x,2),(x+1,5),(x+3,1),(x+4,4)$ & $4$ \\
$Q_3$ & $(1,3,3,1)$ & $(x,4),(x+1,1),(x+3,5),(x+4,2)$ & $4$ \\
$Q_4$ & $(3,3,1,1)$ & $(x,4),(x+1,5),(x+3,1),(x+4,2)$ & $4$ \\
$E$ & $(2,2,2,2)$ & $(x,3),(x+2,1),(x+2,5),(x+4,3)$ & $4$ \\
$L$ & $(1,2,4,2)$ & $(x,3),(x+1,1),(x+4,5),(x+5,3)$ & $5$ \\
$R$ & $(4,2,1,2)$ & $(x,3),(x+1,5),(x+4,1),(x+5,3)$ & $5$ \\
\bottomrule
\end{tabular}
\caption{The twelve templates of tilted rectangles on five rows that touch row $1$, with the
leftmost corner in column $x$.}
\label{tab:fivetemplates}
\end{table}

\emph{Step 2. The six templates that avoid row $1$ are never weak.} Their corners lie in the
rows $2$ to $5$, where a black point sits in a column congruent to $1$ modulo $4$, and their
span is at most $3$, so at most one of their columns is such a column. A column holds two
corners of a rectangle only when $a=c$, which among the six happens only for the unit
diamond, whose two corners in its middle column lie in the rows $r-1$ and $r+1$, of the same
parity. Every double column joins rows of opposite parity, so these two corners are not both
black. Hence each of the six templates has at most one black corner and shows at least three.

\emph{Step 3. The weak translates of the twelve.} We go through Table~\ref{tab:fivetemplates}.
In every case the row-$1$ corner is black unless its column is $9$, $13$, $17$ or $21$, and a
corner in a row other than $1$ is black exactly when its column is in the set $A_j$ of its row.

The unit diamond $D$ is weak exactly when one of its two row-$2$ corners is black, at a column
of $A_2=\{1,9,17\}$, so $x\in\{1,9,17\}$ or $x+2\in\{1,9,17\}$, six cases. The row-$1$
corner $(x+1,1)$ is then black and the other two corners are white. If instead the row-$3$
corner $(x+1,3)$ is black, then $x+1\in\{9,13\}$, the row-$1$ corner $(x+1,1)$ is white, and
three corners show.

Each of $S_1$ to $S_4$ has three corners in the rows $2$ to $4$, in three distinct columns of
span at most $3$, so at most one of them is black. The rectangle is weak exactly when one of
them is black, and then the row-$1$ corner, in a column not congruent to $1$ modulo $4$, is
black too. The black corner can be the row-$2$ corner at $1$, $9$ or $17$, the row-$3$ corner
at $9$ or $13$, or the row-$4$ corner at $5$ or $13$. The row-$4$ corner at $21$ is excluded,
since then another corner would lie in column $22$ or $23$. So each of the four templates has
exactly seven weak translates, listed in Table~\ref{tab:fiveweak}, each showing exactly two
corners in columns not congruent to $1$ modulo $4$.

The templates $Q_1$ to $Q_4$ have two extreme corners in the rows $2$ and $4$ at the
columns $x$ and $x+4$, and two interior corners in the rows $1$ and $5$ at the columns $x+1$
and $x+3$. If an extreme corner is black, then $x$ is congruent to $1$ modulo $4$, and then
the row-$1$ interior corner is black, since its column is congruent to $2$ or $0$ modulo $4$.
For $Q_1$ and $Q_2$ the extremes are $(x,2)$ and $(x+4,4)$, and $x\in A_2=\{1,9,17\}$
holds exactly when $x+4\in A_4=\{5,13,21\}$, so either both extremes are black or neither is.
For $Q_3$ and $Q_4$ the extremes are $(x,4)$ and $(x+4,2)$. Here $x\in A_4$ with
$x+4\le21$ means $x\in\{5,13\}$, and then $x+4\in A_2$, while $x+4\in A_2$ means
$x\in\{-3,5,13\}$. So again both extremes are black, except at $x=-3$, where $(1,2)$ is black
and $(-3,4)$ is white. When both extremes are black, we have $x\in\{1,9,17\}$ for $Q_1$ and
$Q_2$ and $x\in\{5,13\}$ for $Q_3$ and $Q_4$, the five adjacent intervals $[1,5]$ to
$[17,21]$. The row-$5$ interior corner is then white, so each such rectangle shows one point
$(t,5)$ with $t\in\{2,4,\dots,20\}$, ten singletons in all. At $x=-3$ the templates $Q_3$ and
$Q_4$ show two corners, the pairs $\{(-3,4),(0,5)\}$ and $\{(-3,4),(-2,5)\}$. If no extreme
corner is black, the rectangle is weak only if both interior corners are black, the row-$5$
one at a column of $A_5=\{17,21\}$. Column $21$ would put a corner beyond column $21$, so
$x+1=17$ or $x+3=17$, that is $x=16$ or $x=14$, the row-$1$ interior corner at $19$ or $15$
is black, and the pairs $\{(14,2),(18,4)\}$, $\{(14,4),(18,2)\}$, $\{(16,2),(20,4)\}$ and
$\{(16,4),(20,2)\}$ result. So the four templates have $10+2+4=16$ weak translates.

For $E$ a black row-$5$ corner at $x+2\in A_5$ forces $x+2\in\{17,21\}$, so the row-$1$
corner $(x+2,1)$ is white, and both row-$3$ corners are white as well, three visible corners.
So a weak $E$ has a black row-$3$ corner at $9$ or $13$, which gives $x\in\{5,9,13\}$ with the
presentations $\{(5,3),(7,5)\}$, $\{(11,5)\}$ and $\{(15,5),(17,3)\}$.

For $L$ a black corner outside row $1$ needs $x\equiv1$ or $x\equiv0\pmod4$. In the first case
the left row-$3$ corner or the row-$5$ corner is black, which forces $x\in\{9,13,17\}$, and
$x=17$ puts a corner beyond column $21$. In the second case the right row-$3$ corner is black,
which forces $x\in\{4,8\}$, and at $x=8$ the row-$1$ corner $(9,1)$ is white and three corners
show. So $x\in\{4,9,13\}$ with the presentations $\{(4,3),(8,5)\}$, $\{(13,5),(14,3)\}$ and
$\{(18,3)\}$.

For $R$ with $x\equiv1\pmod4$ the left row-$3$ corner would be black at $9$ or $13$, but
then the row-$1$ corner at $13$ or $17$ is white and three corners show. With $x\equiv0\pmod4$
the right row-$3$ corner is black for $x\in\{4,8\}$ and the row-$5$ corner is black for
$x\in\{16,20\}$, and $x=20$ puts a corner beyond column $21$. So $x\in\{4,8,16\}$ with the
presentations $\{(4,3),(5,5)\}$, $\{(8,3),(9,5)\}$ and $\{(16,3),(21,3)\}$.

Adding up, $6+28+16+3+3+3=59$ translates are weak, $12$ of them show one corner and the other
$47$ show two, none is hidden, and the leftmost column used is $-3$.

\emph{Step 4. The presentations are pairwise distinct.} Table~\ref{tab:fivesignatures} lists,
for each group of templates, the rows of the two visible corners and their horizontal
separation. For a weak $S_1$ to $S_4$ with a hidden row-$2$ or row-$3$ corner, the rows of the
two visible corners together with their order and separation identify the template and its
translate, and the same holds for the diamond, the $Q$ templates, $E$, $L$ and $R$, whose
explicit pairs differ in their coordinates. The one possible coincidence is between $S_1$ and
$S_2$, or between $S_3$ and $S_4$, with a hidden row-$4$ corner, where both show their row-$2$
and row-$3$ corners at separation $3$. The two hidden corners would then lie in the adjacent
columns $x+1$ and $x+2$, and both cannot be black, since black points outside row $1$ sit in
columns congruent to $1$ modulo $4$. The twelve singletons are the points $(t,5)$ for even
$t$ from $2$ to $20$, the point $(11,5)$ and the point $(18,3)$, all different. So the $59$
presentations are distinct.

\begin{table}[t]
\centering\small
\begin{tabular}{@{}ll@{}}
\toprule
templates & rows of the visible pair and horizontal separation \\
\midrule
$D$ & rows $2,3$, separation $1$ \\
$S_1$ to $S_4$ & rows $2,3$, separation $3$, or rows $2,4$ or $3,4$, separation $1$ or $2$ \\
$Q_1$ to $Q_4$ & rows $4,5$, separation $1$ or $3$, or rows $2,4$, separation $4$ \\
$E$ & rows $3,5$, separation $2$ \\
$L$, $R$ & rows $3,5$, separation $1$ or $4$, or both in row $3$, separation $5$ \\
\bottomrule
\end{tabular}
\caption{The signatures of the $47$ weak presentations with two visible corners.}
\label{tab:fivesignatures}
\end{table}

\emph{Step 5. No axis-aligned rectangle shows the same points.} An axis-aligned rectangle with
the same presentation as a weak tilted rectangle would have at most two visible corners, and
we check the three kinds that occur. There is no vertical pair and no hidden rectangle in
Table~\ref{tab:fiveweak}.

Each of the twelve singletons $\{(x,r)\}$ has $r\in\{3,5\}$ and sits in an ordinary column
$x$. An axis-aligned rectangle showing exactly $(x,r)$ has three black corners, so column $x$
holds a black corner in a row $s\ne r$, which forces $s=1$, and its other column holds black
corners in the rows $r$ and $1$. So $\{1,r\}\in F$, which is false for $r=3$ and for $r=5$.

The only pair in one row is $\{(16,3),(21,3)\}$. An axis-aligned rectangle through both uses
the columns $16$ and $21$ and a second row black in both, but $B_{16}=\{1\}$ and
$B_{21}=\{4,5\}$ are disjoint.

Every other pair $\{(x,r),(z,s)\}$ has $x\ne z$ and $r\ne s$, and the only axis-aligned
rectangle with both as corners has the further corners $(x,s)$ and $(z,r)$, which would both
have to be black. No visible corner in Table~\ref{tab:fiveweak} lies in row $1$. In all
but five of these pairs both columns are ordinary. The five exceptions are
$\{(5,3),(7,5)\}$ and $\{(15,5),(17,3)\}$ from $E$, $\{(13,5),(14,3)\}$ from $L$, and
$\{(4,3),(5,5)\}$ and $\{(8,3),(9,5)\}$ from $R$, and in each of them the other column is
ordinary. So every such pair has an ordinary column $x$, and the further corner $(x,s)$ in
that column lies in a row $s\ne1$ and is white.

So no axis-aligned rectangle shares a presentation with a weak tilted rectangle.
\end{proof}

\begin{table}[t]
\centering\small
\begin{tabular}{@{}lrl@{}}
\toprule
template & $x$ & presentation \\
\midrule
$D$ & $-1$ & $\{(-1,2),(0,3)\}$ \\
$D$ & $1$ & $\{(2,3),(3,2)\}$ \\
$D$ & $7$ & $\{(7,2),(8,3)\}$ \\
$D$ & $9$ & $\{(10,3),(11,2)\}$ \\
$D$ & $15$ & $\{(15,2),(16,3)\}$ \\
$D$ & $17$ & $\{(18,3),(19,2)\}$ \\
$S_1$ & $1$ & $\{(3,4),(4,3)\}$ \\
$S_1$ & $3$ & $\{(3,2),(6,3)\}$ \\
$S_1$ & $6$ & $\{(6,2),(8,4)\}$ \\
$S_1$ & $9$ & $\{(11,4),(12,3)\}$ \\
$S_1$ & $10$ & $\{(10,2),(12,4)\}$ \\
$S_1$ & $11$ & $\{(11,2),(14,3)\}$ \\
$S_1$ & $17$ & $\{(19,4),(20,3)\}$ \\
$S_2$ & $1$ & $\{(2,4),(4,3)\}$ \\
$S_2$ & $4$ & $\{(4,2),(7,3)\}$ \\
$S_2$ & $6$ & $\{(6,2),(7,4)\}$ \\
$S_2$ & $9$ & $\{(10,4),(12,3)\}$ \\
$S_2$ & $10$ & $\{(10,2),(11,4)\}$ \\
$S_2$ & $12$ & $\{(12,2),(15,3)\}$ \\
$S_2$ & $17$ & $\{(18,4),(20,3)\}$ \\
$S_3$ & $-2$ & $\{(-2,3),(0,4)\}$ \\
$S_3$ & $3$ & $\{(3,3),(6,2)\}$ \\
$S_3$ & $6$ & $\{(6,3),(8,4)\}$ \\
$S_3$ & $9$ & $\{(11,4),(12,2)\}$ \\
$S_3$ & $11$ & $\{(11,3),(14,2)\}$ \\
$S_3$ & $13$ & $\{(15,4),(16,2)\}$ \\
$S_3$ & $14$ & $\{(14,3),(16,4)\}$ \\
$S_4$ & $-2$ & $\{(-2,3),(-1,4)\}$ \\
$S_4$ & $4$ & $\{(4,3),(7,2)\}$ \\
$S_4$ & $6$ & $\{(6,3),(7,4)\}$ \\
\bottomrule
\end{tabular}
\hspace{1em}
\begin{tabular}{@{}lrl@{}}
\toprule
template & $x$ & presentation \\
\midrule
$S_4$ & $9$ & $\{(10,4),(12,2)\}$ \\
$S_4$ & $12$ & $\{(12,3),(15,2)\}$ \\
$S_4$ & $13$ & $\{(14,4),(16,2)\}$ \\
$S_4$ & $14$ & $\{(14,3),(15,4)\}$ \\
$Q_1$ & $1$ & $\{(4,5)\}$ \\
$Q_1$ & $9$ & $\{(12,5)\}$ \\
$Q_1$ & $14$ & $\{(14,2),(18,4)\}$ \\
$Q_1$ & $17$ & $\{(20,5)\}$ \\
$Q_2$ & $1$ & $\{(2,5)\}$ \\
$Q_2$ & $9$ & $\{(10,5)\}$ \\
$Q_2$ & $16$ & $\{(16,2),(20,4)\}$ \\
$Q_2$ & $17$ & $\{(18,5)\}$ \\
$Q_3$ & $-3$ & $\{(-3,4),(0,5)\}$ \\
$Q_3$ & $5$ & $\{(8,5)\}$ \\
$Q_3$ & $13$ & $\{(16,5)\}$ \\
$Q_3$ & $14$ & $\{(14,4),(18,2)\}$ \\
$Q_4$ & $-3$ & $\{(-3,4),(-2,5)\}$ \\
$Q_4$ & $5$ & $\{(6,5)\}$ \\
$Q_4$ & $13$ & $\{(14,5)\}$ \\
$Q_4$ & $16$ & $\{(16,4),(20,2)\}$ \\
$E$ & $5$ & $\{(5,3),(7,5)\}$ \\
$E$ & $9$ & $\{(11,5)\}$ \\
$E$ & $13$ & $\{(15,5),(17,3)\}$ \\
$L$ & $4$ & $\{(4,3),(8,5)\}$ \\
$L$ & $9$ & $\{(13,5),(14,3)\}$ \\
$L$ & $13$ & $\{(18,3)\}$ \\
$R$ & $4$ & $\{(4,3),(5,5)\}$ \\
$R$ & $8$ & $\{(8,3),(9,5)\}$ \\
$R$ & $16$ & $\{(16,3),(21,3)\}$ \\
\bottomrule
\end{tabular}
\caption{The $59$ weak tilted rectangles of the half-strip, each given by its template of
Table~\ref{tab:fivetemplates}, its leftmost column $x$, and its presentation.}
\label{tab:fiveweak}
\end{table}

\begin{proof}[Proof of Proposition~\ref{prop:fivecases}]
Fix $c\le1$ and let $R_1\ne R_2$ be rectangles of the board of columns $c,\dots,21$ with the
same presentation $P$ under $S_c$. If both are axis-aligned, consider the graph $G_{S_c}$. The
six double columns carry six distinct row pairs, so there is no cycle of length four, and a
cycle of length six would be a triangle of row pairs, which is impossible because the six
pairs form a bipartite graph on the rows. The singleton columns are leaves and lie on no cycle.
The board has $22-c\ge21$ columns, so Lemma~\ref{lem:girth} gives $R_1=R_2$. If both are
tilted, then $|P|\le2$ by Lemma~\ref{lem:threecorners}, both are weak, both appear in
Table~\ref{tab:fiveweak}, and Step 4 of Proposition~\ref{prop:fivelist} gives $R_1=R_2$. If one
is tilted and the other axis-aligned, again $|P|\le2$, the tilted one appears in the table, and
Step 5 excludes the axis-aligned one. Every case contradicts $R_1\ne R_2$, so $S_c$ is valid.
The count of $59$, the split $47+12$, the absence of hidden rectangles, the column range
$[-3,21]$ and the two distinctness statements are Proposition~\ref{prop:fivelist}.
\end{proof}
